%%%%%%%%%%%%%%%%%%%%%%%%%%%%%%%%%%%%%%%%
% Compact Laboratory Book
% LaTeX Template
% Version 1.0 (4/6/12)
%
% This template has been downloaded from:
% http://www.LaTeXTemplates.com
%
% Original author:
% Joan Queralt Gil (http://phobos.xtec.cat/jqueralt) using the labbook class by
% Frank Kuster (http://www.ctan.org/tex-archive/macros/latex/contrib/labbook/)
%
% License:
% CC BY-NC-SA 3.0 (http://creativecommons.org/licenses/by-nc-sa/3.0/)
%
% Important note:
% This template requires the labbook.cls file to be in the same directory as the
% .tex file. The labbook.cls file provides the necessary structure to create the
% lab book.
%
% The \lipsum[#] commands throughout this template generate dummy text
% to fill the template out. These commands should all be removed when 
% writing lab book content.
%
% HOW TO USE THIS TEMPLATE 
% Each day in the lab consists of three main things:
%
% 1. LABDAY: The first thing to put is the \labday{} command with a date in 
% curly brackets, this will make a new section showing that you are working
% on a new day.
%
% 2. EXPERIMENT/SUBEXPERIMENT: Next you need to specify what 
% experiment(s) and subexperiment(s) you are working on with a 
% \experiment{} and \subexperiment{} commands with the experiment 
% shorthand in the curly brackets. The experiment shorthand is defined in the 
% 'DEFINITION OF EXPERIMENTS' section below, this means you can 
% say \experiment{pcr} and the actual text written to the PDF will be what 
% you set the 'pcr' experiment to be. If the experiment is a one off, you can 
% just write it in the bracket without creating a shorthand. Note: if you don't 
% want to have an experiment, just leave this out and it won't be printed.
%
% 3. CONTENT: Following the experiment is the content, i.e. what progress 
% you made on the experiment that day.
%
%%%%%%%%%%%%%%%%%%%%%%%%%%%%%%%%%%%%%%%%%

%----------------------------------------------------------------------------------------
%    PACKAGES AND OTHER DOCUMENT CONFIGURATIONS
%----------------------------------------------------------------------------------------                               

% \UseRawInputEncoding

\documentclass[fontsize=11pt, % Document font size
                             paper=letter, % Document paper type
                             %twoside, % Shifts odd pages to the left for easier reading when printed, can be changed to oneside
                             openany, % chapters can start on any page
                             captions=tableheading,
                             index=totoc,
                             hyperref]{labbook}

%\documentclass[idxtotoc,hyperref,openany]{labbook} % 'openany' here removes the
  
\usepackage[bottom=10em]{geometry} % Reduces the whitespace at the bottom of the page so more text can fit

\usepackage[english]{babel} % English language
\usepackage{lipsum} % Used for inserting dummy 'Lorem ipsum' text into the template

\usepackage[utf8]{inputenc} % Uses the utf8 input encoding
\usepackage[T1]{fontenc} % Use 8-bit encoding that has 256 glyphs

\usepackage[osf]{mathpazo} % Palatino as the main font
\linespread{1.05}\selectfont % Palatino needs some extra spacing, here 5% extra
\usepackage[scaled=.88]{beramono} % Bera-Monospace
\usepackage[scaled=.86]{berasans} % Bera Sans-Serif

\usepackage{booktabs,array} % Packages for tables

\usepackage{amsmath} % For typesetting math
\usepackage{graphicx} % Required for including images
\usepackage{etoolbox}
\usepackage[norule]{footmisc} % Removes the horizontal rule from footnotes
\usepackage{lastpage} % Counts the number of pages of the document
\usepackage{float}

\usepackage[ruled, vlined, linesnumbered]{algorithm2e} % For algorithms


\usepackage[dvipsnames]{xcolor}  % Allows the definition of hex colors
\usepackage{epstopdf}
\epstopdfsetup{suffix={}}
\definecolor{titleblue}{rgb}{0.16,0.24,0.64} % Custom color for the title on the title page
\definecolor{linkcolor}{rgb}{0,0,0.42} % Custom color for links - dark blue at the moment

\addtokomafont{title}{\Huge\color{titleblue}} % Titles in custom blue color
\addtokomafont{chapter}{\color{OliveGreen}} % Lab dates in olive green
\addtokomafont{section}{\color{Sepia}} % Sections in sepia
\addtokomafont{pagehead}{\normalfont\sffamily\color{gray}} % Header text in gray and sans serif
\addtokomafont{caption}{\footnotesize\itshape} % Small italic font size for captions
\addtokomafont{captionlabel}{\upshape\bfseries} % Bold for caption labels
\addtokomafont{descriptionlabel}{\rmfamily}
\setcapwidth[c]{10cm} % Center align caption text
\setkomafont{footnote}{\sffamily} % Footnotes in sans serif

\deffootnote[4cm]{4cm}{1em}{\textsuperscript{\thefootnotemark}} % Indent footnotes to line up with text

\DeclareFixedFont{\textcap}{T1}{phv}{bx}{n}{1.5cm} % Font for main title: Helvetica 1.5 cm
\DeclareFixedFont{\textaut}{T1}{phv}{bx}{n}{0.8cm} % Font for author name: Helvetica 0.8 cm

\def\currentYear{2021}

\usepackage{scrhack}
\usepackage[headsepline]{scrlayer-scrpage} % Provides headers and footers configuration
\pagestyle{scrheadings} % Print the headers and footers on all pages
\clearscrheadfoot % Clean old definitions if they exist

\automark[chapter]{chapter}
\ohead{\headmark} % Prints outer header

\setlength{\headheight}{25pt} % Makes the header take up a bit of extra space for aesthetics
\addtokomafont{headsepline}{\color{lightgray}} % Colors the rule under the header light gray

\ofoot[\normalfont\normalcolor{\thepage\ |\  \pageref{LastPage}}]{\normalfont\normalcolor{\thepage\ |\  \pageref{LastPage}}} % Creates an outer footer of: "current page | total pages"

% These lines make it so each new lab day directly follows the previous one i.e. does not start on a new page - comment them out to separate lab days on new pages
\makeatletter
\patchcmd{\addchap}{\if@openright\cleardoublepage\else\clearpage\fi}{\par}{}{}
\makeatother
\renewcommand*{\chapterpagestyle}{scrheadings}

% These lines make it so every figure and equation in the document is numbered consecutively rather than restarting at 1 for each lab day - comment them out to remove this behavior
\usepackage{chngcntr}
\counterwithout{figure}{labday}
\counterwithout{equation}{labday}

% Hyperlink configuration
\usepackage[
    pdfauthor={Hannah Grady, Nicholas Nauman}, % Your name for the author field in the PDF
    pdftitle={Laboratory Journal}, % PDF title
    pdfsubject={labNotebookSeniorProject3}, % PDF subject
    bookmarksopen=true,
    linktocpage=true,
    urlcolor=linkcolor, % Color of URLs
    citecolor=linkcolor, % Color of citations
    linkcolor=linkcolor, % Color of links to other pages/figures
    backref=page,
    pdfpagelabels=true,
    plainpages=false,
    colorlinks=true, % Turn off all coloring by changing this to false
    bookmarks=true,
    pdfview=FitB]{hyperref}

\usepackage[stretch=10]{microtype} % Slightly tweak font spacing for aesthetics

\usepackage[framed,numbered,autolinebreaks,useliterate]{mcode}
\usepackage{todonotes}

% This package is for plotting graphs
\usepackage{pgfplots}

%\setlength\parindent{0pt} % Uncomment to remove all indentation from paragraphs

%----------------------------------------------------------------------------------------
%    DEFINITION OF EXPERIMENTS
%----------------------------------------------------------------------------------------

% Template: \newexperiment{<abbrev>}[<short form>]{<long form>}
% <abbrev> is the reference to use later in the .tex file in \experiment{}, the <short form> is only used in the table of contents and running title - it is optional, <long form> is what is printed in the lab book itself

\newexperiment{example}[Example experiment]{This is an example experiment}
\newexperiment{example2}[Example experiment 2]{This is another example experiment}
\newexperiment{example3}[Example experiment 3]{This is yet another example experiment}

\newsubexperiment{subexp_example}[Example sub-experiment]{This is an example sub-experiment}
\newsubexperiment{subexp_example2}[Example sub-experiment 2]{This is another example sub-experiment}
\newsubexperiment{subexp_example3}[Example sub-experiment 3]{This is yet another example sub-experiment}

%----------------------------------------------------------------------------------------
\newcommand{\HRule}{\rule{\linewidth}{0.5mm}} % Command to make the lines in the title page

\setlength\parindent{0pt} % Removes all indentation from paragraphs

\begin{document}

%----------------------------------------------------------------------------------------
%    TITLE PAGE
%----------------------------------------------------------------------------------------
%\frontmatter % Use Roman numerals for page numbers

%\begin{center}

%

\title{
\begin{center}
\href{http://www.bradley.edu}{\includegraphics[height=0.5in]{figs/logoBU1-Print}}
\vskip10pt
\HRule \\[0.4cm]
{\Huge \bfseries Laboratory Notebook \\[0.5cm] \Large Hardware-in-Loop Plant Modeling}\\[0.4cm] % Degree
\HRule \\[1.5cm]
\end{center}
}
\author{ Hannah Grady, Nicholas Nauman \\ \\\Large hgrady@mail.bradley.edu, nnauman@mail.bradley.edu} % Your name and email address
\date{Beginning September 14, \currentYear} % Beginning date
\maketitle


%\maketitle % Title page

\printindex
\tableofcontents % Table of contents
\newpage % Start lab look on a new page

\begin{addmargin}[0cm]{0cm} % Makes the text width much shorter for a compact look

\pagestyle{scrheadings} % Begin using headers

%----------------------------------------------------------------------------------------
%    LAB BOOK CONTENTS
%----------------------------------------------------------------------------------------

\labday{Wednesday, September 15, \currentYear}
\experiment{Meeting Minutes}
% Place your initial below
HG

Dr. Miah, Mr. Guetz, Mr. Buckner, and I discussed the project design form and work statement. Nick was not present. Dr. Miah, Nick, and myself will be recieving documents and diagrams from AutonomouStuff around September 22nd which will add further detail to our project. We also went over project logistics and discussed using GitHub for project work unless the information became confidential. Project deliverables will be finalized at the next meeting, and weekly meetings will be conducted via Zoom. Nick and I will have at least six hours of lab work each week, not including extra research. Agendas will be sent out before each meeting.

%-------------------------------------------

\labday{Thursday, September 16, \currentYear}
\experiment{Meeting Minutes}
% Place your initial below
HG\\

Dr. Miah, Nick, and I met to recap some of the aspects of the meeting on 9/15 and discuss what we would be doing as students for this project. It is important to bring a notebook to all meetings in order to record the meeting minutes, and that all agendas and meeting minutes should be recorded in the lab notebook. Meetings will be on Thursdays from 4-5pm and Dr. Miah will send out meeting invitations to everyone. Lab work will take place on Tuesdays and Thursdays from 4-7pm. Dr. Miah will join these lab sessions either in-person or virtually. We will be using camel case convention for the naming of all folders and files, and Dr. Miah will provide the templates we need. Nick and I will send our GitHub information to Dr. Miah so we can access the senior project GitHub, and Dr. Miah will share a Google Drive folder with us. We discussed the format for lab notebook entries, and Nick and I will need to install Inkscape and TeX Live. 

% -------------------------------------------

\labday{Thursday, September 23, \currentYear}
\experiment{Meeting Minutes}
% Place your initial below
HG\\

Dr. Miah, Mr. Guetz, Mr. Buckner, Nick, and I met to further discuss the project. Mr. Guetz provided block diagrams for the systems that we will be modeling. Everything highlighted in red is related to the vehicle, while everything in blue is related to firmware. The torque voltages range from 1.8-3.8V and based on the voltage the steering wheel will turn either left or right. AutonomouStuff already has a steering controller, but it has trouble with small steering angles. Some deliverables we discussed were having a Simulink model that mathematically aligns with data along with a report, test plans for data collection, small steering angle and mathematical model, and results. Once we get raw data from an actual vehicle we will then come up with a Simulink model. A tentative schedule we discussed was spending the first semester gathering data and modeling that data, and then conducting experiments during the second semester. We might go out to AutonomouStuff to collect data, right now we need to start thinking about what type of data we need to collect and formulate a test plan. When we do this depends on vehicle availability, and AutonomouStuff will let us know what days the vehicle will be available. Dr. Miah, Nick, and myself were granted access to the BitBucket repository for this senior project by AutonomouStuff. This is where all documents from this meeting will go. At the end of the meeting Nick asked if there were any ranges we needed to focus on for data collection. The answer is that we need to look at speeds, with 20mph being the max limit. 

% -------------------------------------------
%-------------------------------------------

\experiment{Background study}

HG\\
During the lab, I worked on reading and taking notes of the different sections of the "Get Started with Sytem Identification Toolbox" documentation in MathWorks. I was able to read most of the documentation in the "About System Identification" section, but I will have to finish reading the rest of it and completing the tutorials over the weekend. 

\vspace*{12pt}

NN\\
Today during the lab time, I researched how the Mathworks System Identification Toolbox allows users to create mathematical models of dynamic systems from measured input and output data using MATLAB functions and Simulink in either time-domain or frequency-domain. The research was completed using the following documentation: \href{https://www.mathworks.com/help/ident/getting-started-1.html} \\
There are multiple methods for obtaining identification, but I believe the most effective method for this project will be to capture important system dynamics and organize the data into variables. Once the data is assigned to variables, it can be imported into the System Identification app or set as a data object that can be used to estimate a model from the command line. The System Identification Toolbox software supports time-domain and frequency-domain data for creating linear models and only time-domain data for creating nonlinear models. The data can be single or multiple inputs and outputs that are either real or complex. Once the data is prepared and ran through the System Identification Toolbox, there are functions that are used for either Linear Model Estimation, Nonlinear Model Estimation, or Grey-Box Model Estimation methods. Once the model equations are built, there are options for model validation (comparing experimental data with expected model data), Model Analysis (Model Transformation, Data Extraction, Simulation and Prediction, and Response Computation), and Time Series Analysis.

% -------------------------------------------


\labday{Tuesday, September 28, \currentYear}
\experiment{Background Study}
% Place your initial below
HG\\
During part of the lab time we met with Dr. Miah. We discussed how a folder called "labnotebook" was added to BitBucket. The Virtual Desktop Matlab application doesn't have the System Identification toolbox, Dr. Miah will ask about getting it added to the virtual desktop. We talked about what we have been working on in lab so far. Today I mainly read through some System Identification Toolbox documentation and found articles for the literature review. 

NN\\
During the lab time after meeting with Dr. Miah, I worked on reading documentation on the Matlab System Identification Toolbox and looking at examples of using the System Identification Toolbox. Also, I started creating the block diagrams for the vehicle subsystems that we will be collecting data for and modeling. I got the steering system, acceleration system, shift system, and speed control system block diagrams created. 

% -------------------------------------------

\labday{Thursday, September 30, \currentYear}
\experiment{Meeting Minutes}
% Place your initial below
Agenda:
\begin{itemize}
	\item Discuss a plan for data collection of the Lexus car platform
	\item Discuss the logistics of when data collection may occur
\end{itemize}

NN\\
In today's meeting, Dr. Miah, Mr. Guetz, Hannah, and I discussed the ranges of the inputs and outputs of the steering model so a data collection test plan could be created. Also, we discussed the logistics of collecting this data. The Lexus platform will not be available for testing for at least two weeks. This should allow Mr. Guetz enough time to get the firmware ready for collecting data.

\experiment{Background study}

HG\\
Today during the lab time I started reading through documents that could be used in the literature review section of our proposal and summarizing them. I also watched some videos on the MathWorks website about using the System Identification Toolbox. 


%-------------------------------------------

\labday{Tuesday, October 5, \currentYear}
\experiment{Background Study}
% Place your initial below
HG\\
During the lab time Nick and I started working on the System Level Functional
Requirements presentation. I added to the Problem Statement, Literature Review,
and References slides. We also met with Dr. Miah to go over the expectations for
the presentation and so he could give us feedback on our draft of the System
Level Functional Requirements document. Dr. Miah also pushed the changes we made
to some of the documents to the GitHub repository we have for this project.

NN\\
During the lab time, I added the Introduction, System Inputs/Outputs, System Requirements, and Concluding Remarks sections of the System Level Functional Requirements presentation. 

%--------------------------------------------------------------------------------------------

\labday{Thursday, October 7, \currentYear}
\experiment{Data Collection}
% Place your initial below
HG\\
Today Dr. Miah, Nick, and I went to AutonomouStuff in Morton to collect data. We met Mr. Guetz and Mr. Buckner there, and they helped us sign in and showed us where the vehicles were kept. Dr. Miah, Nick, Mr. Guetz, and I went out to the Lexus and used that for testing. Dr. Miah was the safety driver while Mr. Guetz used a computer to give commands to the vehicle. Nick and I took notes. There is a red button to shut down the system, and the safety driver can also turn the steering wheel or step on the gas in case anything goes wrong. We started measuring data for the steering model first. At rest the steering angle is zero, and if the driver tries to manually move the steering wheel while the software is engaged, the torque voltages will fight them. Turning the wheel counterclockwise corresponds to positive radians, while clockwise corresponds to negative radians. The wheel angle is proportional to the steering angle. Dr. Miah turned the steering wheel all the way to lock and back twice to collect that data. We then moved onto collecting data at 0.1 radians to the left and then to the right multiple times. From there we moved onto pi/2 radian turns at 5mph. We also collected data for multiple steering angles at 10mph. We then kept collecting data, keeping the car at a speed between 10-11mph. Mr. Guetz explained that the car can also be controlled using a joystick, but that tends to be jittery. Data can also be collected while the car is being driven manually, and we ended up collecting data while the car was being controlled manually and automatically. From the computer, Mr. Guetz requests a position, and the car produces a single torque voltage that gets split into two torque voltages. Torque is the amount of assist applied to the rack. Once we had collected the steering data, we then moved onto collecting brake data. We collected that data while parked, and gathered data while the car was being operated manually and automatically. In automatic mode, there is no visible pedal change, since everything happens internally. We did two rounds of testing, during the first round we collected data manually at a range of 0-93 and then automatically from 0-100. For the second round we tested the brakes manually from 0-98 and automatically from 0-100. We got the data from this in terms of voltages. At the range of 0-75, the brake pressure is low, and you reach a maximum pressure as you approach 100. Afterwards we collected data for our acceleration model. It was similar to the data collection for the brake model in that the vehicle was put in park, only this time the acceleration pedal was pushed. We did two rounds of data collection in manual and automatic. The first round tested the acceleration from 0-100 for the vehicle in manual and automatic, and did the same thing for the second round. The input to the system is the voltages and the output is the pedal position. We will come back to AutonomouStuff later to collect data for the shift and gear models. 

NN\\
After completing the testing as described above, we returned to the garage to ask any remaining questions. I asked Mr. Guetz about the hardware setup for each system we are developing models for. He explained that the hardware implementation each similar for each system. The system is connected to the PACMod which is then connected to a CAN case, which connects to the laptop. The laptop can be used to log this CAN data, which is how we collected the data we will be analyzing. 

%--------------------------------------------------------------------------------------------
\labday{Thursday, October 14, \currentYear}
\experiment{Background Study}
% Place your initial below
HG\\
At the beginning of the lab time, Nick and I met with Dr. Miah. He let us know that he had uploaded the pictures from our trip to AutonomouStuff. Nick mentioned that he had started using the System Identification Toolbox with the data we collected last Thursday. We went over a block diagram for the steering model. There are two torque voltages A, and B, so the model will have to be updated. The controller deals with the error between the desired steering angle (which Eric was setting during data collection) and the actual steering angle. We discussed how we need to understand the physical significance of A and B and what they do. We went over the data that Nick already had open in Matlab. The commanded value might not be a percentage, and instead may be in radians. We will have to ask Eric about this. The steering RPT Output value is the actual steering angle. The actual steering wheel angle is based on output voltages A and B, and using these values will help us to identify our two-input two-output system. Nick and I will have to do some research on what order system reliably gives a good model. We will also have to do research and see how people come up with different models based on a two-input, one output system. We will include any research in our proposal. Dr. Miah will ask Mr. Mattus about having the System Identification Toolbox on the lab computers. 

NN\\
During the lab time, Dr. Miah, Hannah, and I discussed the data we collected during our trip to AutonomouStuff. Dr. Miah showed us how to develop a block diagram to describe the behavior of the steering model. This allowed us to easily select which data we needed to use as the inputs and outputs for using in Matlab. With these details, I started working on trying to develop a model of the steering motor using the data collected from AutonomouStuff. I ran into some problems with not add my data correctly so I researched some articles that explained how to process the data.

%-------------------------------------------

\labday{Tuesday, October 19, \currentYear}
\experiment{Background Study}
% Place your initial below
HG\\
At the beginning of lab Nick and I met with Dr. Miah to discuss what we have been working on. We talked about the steering model and how the system had two inputs (torque voltages A and B) and one output (Actual SW angle). We discussed the possible options for identifying a model for the steering system using the System Identification Toolbox. One option is to use A and B as  the inputs and the commanded (desired) steering angle as the output. We talked about creating an agenda for Thursday's meeting that includes questions we have for Erik. This agenda will be placed in BitBucket along with our modified diagrams. The System Level Functional Requirements Document must be finalized by October 21st, along with the presentation. Nick and I will work on this outside of lab time. During the lab time Dr. Miah wanted us to work with the data and find examples showing system modeling of MISO systems using the toolbox. Nick has already started working with the Data before our meeting today. He has generated a transfer function, and wants to tune the generated parameters. We will try all possible options and see what gives the best fit and validation data for a model. Since we will have two inputs, when both inputs are activated, the ultimate output is the superimposed result. The inputs should be coupled, so we will need to modify the way we are simulating. 

During the lab time I worked on finding examples of system identification of MISO systems. I was able to find a couple related documents, including one Matlab tutorial. I worked through this tutorial during the lab time. 
\begin{itemize}
	\item \href{https://www.mathworks.com/help/ident/gs/identify-linear-models-using-the-command-line.html}{Example}
\end{itemize}

NN\\
During the lab time, Dr. Miah, Hannah, and I discussed the data that we collected from AutonomouStuff. We discussed the need to get clarification from Erik and Joe about what units were being used for the actual steering angle, commanded steering angle, and reference steering angle. We also discussed getting clarification on what the non-linearities actual mean to the system. After the discussion, I worked on manipulation the data in the System Identification Toolbox to try and get a model that closely follows the system output. 

%---------------------------------------------------------

\labday{Thursday, October 21, \currentYear}
\experiment{Meeting Minutes}
% Place your initial below
HG\\
I arrived late to today's meeting with the AutonomouStuff since my class doesn't get out until 4:15pm. When I entered the meeting Nick was sharing his screen and the data that he had so far. He mentioned that the parameters needed to be tuned up so that the desired steering wheel angle matched the actual steering wheel angle. The current model he had wasn't very good at predicting the steering angle but it got better as time went on. We will have to try tuning a model to handle small steering wheel angles and then apply that model to the larger steering wheel angle data to see if it works. Dr. Miah's plan is to create data sets out of the data that we have and then generate random test sets. We will then isolate some data sets where the response is small steering wheel angles and make test data. Dr. Miah drew a new steering wheel model block diagram and then two options. The first option was to use the output torque voltages A and B from the steering wheel model as inputs to a predictive steering model, with the output being the actual steering wheel angle. The second option was to take the torque voltages A and B from the controller output as our predictive model's inputs and then have the output be the commanded steering wheel angle. The problem is mainly coming up with a model and then tuning the parameters in such a way so the real data matches with the predictive data. After getting a model, it will be tested with AutonomouStuff's controller to see if they can get better control. Our model will have to have the non-linearities seen in the real world. Mr. Guetz mentioned that a small steering wheel angle is considered to be between -10 and 10 degrees. We also talked about how a small steering angle means that from time to time the steering angle varies very little. Mr. Guetz explained that the non-linearities are larger in smaller angles, especially when the car is trying to stay in one lane. These small non-linearities cause the car to swerve back and forth. He then drew us a diagram to help illustrate the non-linearities, which occur when going to the center. We will need to add a section on non-linearities to our report. We will need to predict our model in such a way that it will respond to small non-linear angles faster than the existing system. As we go along we will explore a deep reinforcement learning algorithm. Erik will run through the data and make sure that it is okay. He will then send it to us by Monday. For now we will work with the clean data and then work with the rest of the data.

\experiment{Background Study}
% Place your initial below
HG\\ 
After our meeting with AutonomouStuff, Nick, Dr. Miah, and I met. Dr. Miah told us to write everything and place all diagrams and plots in the lab notebook and project proposal. Going forward we will be working on the proposal template. We will start using GitHub for all document templates so Nick and I can edit simultaneously. Nick and I divided up the sections of the proposal to work on and I looked for examples of using system identification for MISO or MIMO systems. 

NN\\
Dr. Miah, Hannah, and I met with Joe and Erik to discuss clarifications with the units of the steering angle since there was some confusion from an email Erik sent. We also got some clarification on what the system non-linearities were and how the affected the steering. Erik also provided some context to how the models would be used once we completed them. While we wait for Erik to send an email with what is the bad data in the data sets, Dr. Miah advised us to do more research on using the System Identification Toolbox for modeling MISO systems and examples we could run through. These are some of the articles I found:
\begin{itemize}
	\item \href{https://www.mathworks.com/help/ident/ug/dealing-with-multi-variable-systems-identification-and-analysis.html}{Example 1}
	\item \href{https://www.mathworks.com/help/ident/ug/modal-analysis-of-a-flexible-flying-wing-aircraft.html}{Example 2}
	\item \href{https://www.mathworks.com/help/ident/gs/system-identification-workflow.html}{Workflow}
	\item \href{https://www.mathworks.com/help/ident/model-validation-basics.html}{Validation Basics}
\end{itemize}


%-------------------------------------------

\labday{Tuesday, October 26, \currentYear}
\experiment{System Identification}
% Place your initial below
HG\\
Dr. Miah, Nick, and I met at the beginning of the lab time to discuss our progress. I had read a paper over the weekend that mentioned modeling a MISO system, and so during the lab time I finished summarizing it and placed that summary in the proposal document. There is an error that keeps popping up that I tried to resolve, but after doing some research, I still wasn't able to fix the problem. I will have to email Dr. Miah about this. It did not affect my ability to add stuff to the proposal. I also wrote a mathematical equation for the summary and also placed that in the project proposal. I then found another example in Matlab about dealing with MIMO systems and estimating models, so I worked on that during the lab time. 
\begin{itemize}
	\item \href{https://www.mathworks.com/help/ident/ug/dealing-with-multi-variable-systems-identification-and-analysis.html}{Example}
\end{itemize}

NN\\
During the lab time today, Dr. Miah advised Hannah and I to work through some examples for using the Matlab System Identification Toolbox, specifically for MISO or MIMO systems so that we can verify we are working through the process correctly. I worked through the following examples and when I was confused by what a certain function did, I looked at the documentation to understand it better. 
\begin{itemize}
	\item \href{https://www.mathworks.com/help/ident/ug/dealing-with-multi-variable-systems-identification-and-analysis.html}{Example 1}
	\item \href{http://www.ece.northwestern.edu/local-apps/matlabhelp/toolbox/ident/ch3tut44.html}{Example 2}
\end{itemize}

Also, Dr. Miah announced that the Matlab System Identification Toolbox was installed on computers 9 and 10 in the Senior Lab.

\labday{Thursday, October 28, \currentYear}
\experiment{Meeting Minutes}
% Place your initial below
NN\\
During the lab time, Dr. Miah, Erik, Hannah, and I met to discuss getting clean data to work with for the steering data we collected. Erik said he could get us clean data by Monday, and I would send a reminder email if he forgot to upload the data. After the meeting, Hannah and I worked on summarizing the Matlab System Identification Toolbox examples we found into the proposal document. After that, we worked on trying to use the System Identification Toolbox on the steering data we collected.

\experiment{Background Study}
% Place your initial below
HG\\
After our meeting with AutonomouStuff, Nick, Dr. Miah, and I met to discuss what we would be working on for the rest of the lab time. Dr. Miah reminded us to not export any whitespace from Inkscape. We also talked about adding the matlab examples we worked on to the lab notebook. From now on, we will be pulling the repository from GitHub before starting work on any documents. During the lab time I worked with the data that we had in Matlab and also started summarizing the matlab examples to be placed in the project proposal. 

NN\\


\labday{Tuesday, November 2, \currentYear}
\experiment{Lab Time}
% Place your initial below
HG\\
Dr. Miah, Nick, and I went over the things to be done at the beginning of the lab time. Dr. Miah showed us how to create references using BibTex. We also talked about the draft proposal that we were submitting that day. We had received the data from Erik, and so we started working on creating models using the clean data during the rest of the lab time. 

NN\\
During the lab time, Dr. Miah, Hannah, and I met to discuss what items should be worked on. There is a draft proposal document due by the end of today that Dr. Miah reviewed and advised us to submit. For the rest of the lab time, Hannah and I worked on developing models of the steering system in the Matlab System Identification Toolbox so the models could be reviewed with AutonomouStuff on Thursday.

%-------------------------

\labday{Thursday, November 4, \currentYear}
\experiment{Meeting Minutes}
% Place your initial below
HG\\
At the beginning of lab Dr. Miah, Erik, Nick, and I met to talk about our progress. Going from the meeting I worked on modeling the manual and by-wire steering angle systems. 

NN\\
During the lab time, Dr. Miah, Erik, Hannah, and I met to discuss the models we developed to get feedback. Hannah presented her state-space model that had a best fit percentage of 59. I presented my transfer function model that had a best fit percentage of 96. Erik stated that these models still did not track close enough to the non-linearities, but were accurate almost every else. We talked about each of the logs. Logs 1 and 3 were manual steering angle data, while Log 2 was a commanded steering angle (By-Wire method). Erik recommended running two different training sessions: one for the By-Wire method and one for the manual method. He recommended this because for the By-Wire method, the electronic control unit (ECU) gets rid of the input torque voltages by open circuiting the motors and instead sending output torque voltages from the controller to the motors. For the rest of the lab time, I worked on running the separate training sessions for the steering system model and then started preparing data for the acceleration model.

%-------------------------

\labday{Tuesday, November 9, \currentYear}
\experiment{Lab Time}
% Place your initial below
HG\\
At the beginning of the lab time Dr. Miah and I met. We briefly went over the proposal presentation and he mentioned that we have to have all of the figures in the presentation in the proposal document. He will work on getting the Gantt chart template added to the presentation so we can use it in the Timeline and Milestones section. During the lab today he asked me to work on getting the raw data and plotting the actual and desired steering wheel angles along with the measured and controller A and B voltage values. Afterwards I need to place these plots into the proposal document. He also corrected some of the errors that were showing up in the proposal document. During lab I worked on generating these plots, and I also started working on modeling the acceleration subsystem. I have some questions about modeling it so I didn't get as far as I had hoped. 

NN\\
I was unable to attend the meeting during lab time due to a conflict. I worked on modeling the acceleration subsystem on Wednesday, November 10. I was able to derive models that are very accurate for both by-wire and manual methods. I have attached the results in the proposal presentation.

%-------------------------

\labday{Thursday, November 11, \currentYear}
\experiment{Meeting Minutes}
% Place your initial below
HG\\
I arrived late to the meeting due to a doctor's appointment. Dr. Miah, Erik, Joe and I met at the beginning of the lab time to go over the presentation and the steering system model. During the lab time I worked on developing a model for the brake subsystem.  

NN\\
Dr. Miah, Erik, Joe, Hannah, and I met to discuss the models we had developed so far. I presented the models developed for the steering system and acceleration system for discussion and feedback. Dr. Miah and Joe recommended adding more labels to the graph for giving a clear message of what the plot means. We also looked at the error plots to determine if the error was within an acceptable range and to determine the behavior (quick spikes or smooth plot). Also, we discussed implementing our developed models on the AutonomouStuff HIL bench to determine the accuracy of the model with new data. Erik suggested adding the developed model into the HIL bench and running it in parallel with the current model AutonomouStuff uses and we can compare the effectiveness during the model. For the lab time, I worked on adding labels and other details to the figures of the developed models of the steering and acceleration system. I also added the mathematical equations for these models.

%-------------------------

\labday{Tuesday, November 16, \currentYear}
\experiment{Lab Time}
% Place your initial below
HG\\
At the beginning of lab Dr. Miah, Nick, and I went over the proposal presentation. He offered suggestions on what to change, such as making the graphs larger and adding a timeline to the timeline and milestones slide. He also mentioned that all changes we make and any pictures we add should also be added to the proposal document. During the lab time I worked on the presentation by adding in a MATLAB example for the Preliminary Work slide and adding a break to the references slide. I also made some small changes to the title and concluding remarks slides. I then added a timeline for the timeline and milestones slide. Nick and I also divided up what slides we would present on Thursday. 

NN\\
After reviewing the suggestions Dr. Miah provided, I worked on adding figures for the simulation results, adding tables for the coefficients of the model transfer functions, adding an experimental setup picture to that slide, and making some edits to wording on some of the slides I would be presenting.

%-------------------------

\labday{Thursday, November 18, \currentYear}
\experiment{Meeting Minutes}
% Place your initial below
HG\\
At the beginning of lab, Nick and I gave our proposal presentation to Dr. Miah, Erik, and Joe. Once we were done presenting they gave us feedback on how we could improve our presentation. During the lab time I worked on editing the parts list, concluding remarks, and problem statement slides. 

NN\\
After Hannah and I gave our proposal presentation to Dr. Miah, Erik, and Joe and they provided feedback, the rest of the lab time was spent updating the slide deck based on the feedback. I fixed the figures of the simulation results and the wording of the experimental setup. After that, I worked on adding some of these details to the proposal document.

%-------------------------

\labday{Tuesday, November 23, \currentYear}
\experiment{Lab Time}
% Place your initial below
HG\\
During the lab time I worked on making the font for the inputs and outputs of the steering, brake, and speed model block diagrams. The acceleration, shift, and speed control model block diagrams didn't show up when I opened them in Inkscape. I then finished reading through the proposal document and made some small changes to it. 

NN\\
During the lab time I worked on plotting the graphs of the data for each vehicle subsystem compared with the model we developed to compare the accuracy of the actual data with the estimated plot. I added units and labels to all of the figures to ensure the graph was clear. Once I had the figures created, I added them to the proposal document.

%-------------------------

\labday{Tuesday, November 30, \currentYear}
\experiment{Lab Time}
% Place your initial below
HG\\
Dr. Miah and Nick were working on proposal edits when I joined the Zoom call. Dr. Miah pushed the changes made to the repository. We went over some changes Nick and I have to make to the proposal and we talked about changes for the proposal presentation. We decided that we would work on our proposal presentation on Thursday at 3:30pm and do a mock presentation tomorrow at 4pm. I sent out a calendar invite for this mock presentation. During the lab time I worked on making changes to the proposal and proposal presentation. This included updating fonts for the block diagrams of the shift, acceleration, and speed control models. 

NN\\ 
Dr. Miah, Hannah, and I worked on polishing the proposal with edits to the figures of our system identification results, as well as adding details to the introduction and background sections. I also added some of the details to the proposal presentation as well.


\labday{Monday, January 24, \currentYear}
\experiment{Meeting Minutes}
% Place your initial below
HG\\
Dr. Miah, Nick, and I met to discuss the upcoming semester. Dr. Miah stressed that in ECE 499 we will focus more on implementing our senior project using computers, along with testing/validating our models using HIL. We decided that the lab times for this semester will be on Mondays and Wednesdays from 2-5pm. Nick will send the zoom invite for the lab times. In addition to lab, there will also be a weekly meeting with representatives from AutonomouStuff. We are still working on setting up a time for that meeting, and I will send an email to all involved detailing my availability and Nick's availability. During the meeting we also discussed the deliverables for this semester, and Dr. Miah asked us to make changes to the proposal document reflecting these deliverables. The last deliverable of the semester will be the final report, and Dr. Miah will provide an IEEE template for us to use. He also reminded us to make sure to take a lot of pictures during our next trip to AutonomouStuff to include in the report. 

\labday{Wednesday, January 26, \currentYear}
\experiment{Lab Time}
% Place your initial below
HG\\
Nick and I met in the senior lab and worked on writing paragraphs for the proposal document about each deliverable for this semester. We compared models for the steering and acceleration subsystems to make sure that we were on the same page. I showed him what I had so far for the brake subsystem. I will work on modifying the code I have for the brake subsystem to better match the format he has been using for the steering and acceleration subsystems.  

NN\\ 
Hannah and I met in the senior lab so we could compare the steering, acceleration, and brake models that we had developed. I also worked on creating a script that simplified adding the data to the system identification tool box GUI and then I worked on refining the steering system model.


\labday{Monday, January 31, \currentYear}
\experiment{Lab Time}
%Place your initial below
NN\\
During the lab, I started analyzing the effects of adding/removing zeros and poles from the transfer function model that I developed for the steering model. Further analysis will need to be completed to determine if the transfer function models meet the accuracy and complexity requirements we discussed.

HG\\
During the lab I kept working on the brake model. I had some questions for Nick about his formatting and he showed me how he was importing logs for the data. I realized that my assumptions of which brake logs were manual and which were by-wire was incorrect, so I went back through the logs and determined which were manual and which were by-wire. 

\experiment{Meeting Minutes}
HG\\
Near the end of the lab time Dr. Miah, Erik, Nick, and I met to discuss our progress and what would be accomplished during the upcoming semester. Nick updated the group on the progress we had made so far. We currently have models created for the steering, acceleration, and brake subsystems created. We are close to starting testing and we also need to collect data in order to model the remaining subsystems. Nick asked Erik if we needed to find models using by-wire data since we are working to develop a new controller. Erik suggested that we should focus on using only manual data to create our models from now on. Dr. Miah mentioned that we will write a section in the report detailing why using the by-wire data doesn't make sense. Nick also asked about computing power considerations, since we do not want to come up with a model that needed more computational power than the system could provide. Erik said that the model will be put on a machine that can handle complex systems, so computational power isn't an issue. The thing to keep in mind with a complex model is that it is harder to fine tune lead and lag. After Nick asked about the desired accuracy threshold, Erik said that it should be within 0.1 radians. We decided that using the manual models, we will increase the number of poles and zeros and see how the model changes. We will also find the locations for the poles and zeros. It was said that this analysis should lead to accuracy within the accepted thresholds. Dr. Miah then asked Erik about HIL testing availability in February. Erik said that the earliest we could do this testing would be in late February, and mentioned that the models would have to be MATLAB readable. AutonomoStuff uses Simulink, so we will probably make our models into Simulink blocks. AutonomouStuff will provide the virtual environment and the inputs for testing, and Erik can give us a screenshot of where our models should go in Simulink that we can use for our report. We concluded that we would plan on conducting testing during the last week of February and that Erik would let us know what dates would work. We will also get back to Erik on what dates would work for us and adapt as necessary with a second visit in April or the end of March to ensure that the models are working. Our action items for this week are to work on the models and analyzing the transfer functions by adding poles and zeros along with creating a plan and time for subsystem data collection. We will also finish making changes to the proposal document and upload it to BitBucket. Dr. Miah is going to commit his changes to the document as well. 


\labday{Wednesday, February 2, \currentYear}
\experiment{Lab Time}
% Place your initial below
NN\\ 
During lab time, I worked on adding a section to the proposal document outlining our intention to only use the manual mode data for our models. I also added an annotated figure of how the steering system works. After that, I continued my analysis of the steering system transfer function model.

HG\\
During the lab, I worked on making changes to the spring 2022 Gantt chart and rearranging the paragraphs detailing our semester deliverables. I also added some more headings for these paragraphs. Next, I added my meeting minutes from Tuesday to the lab notebook and pushed those changes. I went over the changes Nick and I made in the proposal and added a few more changes. I then kept working on the brake model. 

\labday{Monday, February 7, \currentYear}
\experiment{Lab Time}
% Place your initial below
HG\\ 
During the lab time I worked on analysis of the brake subsystem using different transfer functions. I removed the operatorInteraction output from my brakeLog model and removed the mean values from them as well, to keep it consistent with what Nick has been doing with the steering and accel models. Using the manual part of log 2 to create the model and the manual part of log 3 to validate it, I found that I got a best fit percentage of 97\%
using a system that has 20 poles and 2 zeros. The second best system had a best fit percentage of 96.97\%
and had 21 poles and 2 zeros. I started analyzing using brake data from log 1 as well. I split up the log 1 data in half, using the second half to create the system and the first half to validate it. I found that a transfer function with 20 poles and 2 zeros had a best fit percentage of 71.85\% 
and the system with the best fit had two poles and one zero. The two poles and one zero transfer function gave a best fit percentage of 86.1\%
. 

NN\\

During the lab time, I worked on analysis of the steering subsystem transfer function models. As part of the analysis, I created a script that would  plot and compare transfer functions for poles up to 30 using each log of data as both the estimation data and the validation data with respective runs. The script was not completely functional at the end of lab time, but it should be working for the next lab period.


\experiment{Meeting Minutes} 
HG\\ 
Near the end of lab Dr. Miah, Nick, Erik, and I met to go over our progress and what our next steps will be. I presented what I had found during brake system analysis. We discussed whether we needed to find the relationship between on/off and brake pressed, since that relationship is already known. I said that I would write down the transfer function that gives the best fit percentage. Once we make the transfer function block, Erik said that we will be able to put it in with the HIL bench. We will be using the model for testing and not the actual brake subsystem. A question about checking the stability of the system on the car was asked. Erik said that he wasn't sure if there was a metric for that, but that the biggest thing to look for is if there is a large tracking error. We decided that we will test the system stability using the step response in Simulink. We will do this with each subsystem. This will help us to determine whether our systems are stable before we do HIL testing. Nick asked about when we would be able to test our systems using HIL. Erik said that he is waiting from a response from the HIL engineer, and that Joe said we are okay on vehicle availability. He will get back to us on what the HIL engineer says. Nick mentioned that we are still doing analysis but once we are done we can create transfer function blocks in Simulink. Dr. Miah let Erik know that we will send the sim files to him about a week or two before we come in for testing to verify that they are compatible with the system. We will need to be able to just drag and drop the Simulink files into the HIL system. This should be done by the end of next week. Erik will work on getting a picture or diagram of how everything on the HIL testbench will look. This will help us to see the inputs and outputs our Simulink blocks will need to account for. Dr. Miah mentioned that we will put everything we have so far into the proposal document and will update it in BitBucket and GitHub by Wednesday. We will have to create a section about the brake system and we will need to update some block diagrams. Near the end of the meeting Erik mentioned that AutonomouStuff has never really done any modeling before, and this is the first time that they will have a model of the vehicle. By using the HIL bench and just writing code, it saves time and money that would have been spent on vehicle upkeep. Dr. Miah let Erik know that we will write about this in the proposal introduction and abstract. Our action items for this week will then be: 
\begin{itemize}
\item 1. Finish analysis of brake, accel, and steering subsystems
\item 2. Update the proposal and upload it to GitHub and BitBucket
\item 3. Create Simulink blocks
\item 4. Test systems using step response
\item 5. Send sim files to Erik
\end{itemize}
   
\labday{Wednesday, February 9, \currentYear}
\experiment{Lab Time}
% Place your initial below
NN\\
During the lab time, I finished the script that would develop a model using each log of data as the estimation data and then compare that model with all the other logs as the validation data for a varying number of poles. After running this script, I could not find any model that fit our accuracy and error bound requirements. Dr. Miah advised me to start looking into using neural networks for developing the models.

HG\\
During lab I continued analysis of the brake model, trying to find the transfer function that would give the best fit. So far the model with 28 poles and 10 zeros gives a best fit within the acceptable range of accuracy for most of the log combinations, but not for all. So far I have been unable to find one that gives a best fit percentage that is within the acceptable range of accuracy for all of the log combinations. 

\labday{Monday, February 14, \currentYear}
\experiment{Lab Time}
% Place your initial below
NN\\
During the lab time, I started researching into how to use MATLAB's deep learning toolboxes to develop the system models. Dr. Miah advised me to use the Neural Network Time Series app to develop the model. I worked through some examples online and then I started using the data we collected to develop a model of the steering system. We did not meet with Erik today as he had other commitments. He did tell us that we would be able to test our models on a HIL bench as early as late next week. We emailed him to let him know our availability on Friday in case the HIL bench was ready by then.

HG\\
During lab I started reading through the documentation Dr. Miah sent about the Deep Learning toolbox in MATLAB. I started looking at the error plots for the brake models. Nick emailed Erik about what the error bounds for the brake and acceleration models would be, but for now I just assumed a bound of 5\%.  

\experiment{Meeting Minutes}

HG\\
Near the end of the lab time Nick, Dr. Miah, and I met to discuss some questions Nick and I had. Dr. Miah had Nick google nn to svg, then explained that training means that we are trying to find the optimal weights for each connection. We need to come up with an architecture and then go to tools that are already available. Once we have a system we should be able to input any data and still get an accurate output. We need to find the mean and variance of a data set in MATLAB, and then we can use the statistics toolbox. Once we do that we should try and see if the values for the different data sets are the same, but only by looking at some of the data (ex. torque voltages). Dr. Miah suggested that we start with the Neural Net Fitting and then the Neural Net Time Series toolbox. He mentioned that we should read through the documentation and write down the theory behind it along with writing down all of the steps. Dr. Miah will be able to give us feedback on this. He also mentioned that he converted the proposal into the final report format, and that from now on we should only be editing the report. Nick and I will have the acceleration and brake models reading for testing. We decided that we will go to AutonomouStuff for data collection and model testing either next Thursday or Friday. After the meeting Nick emailed Erik asking about the error bounds for the brake and acceleration models, if he had a skeleton diagram of how the HIL bench would be set up, and if we could come to AutonomouStuff next Friday afternoon for testing and to collect data. 


\labday{Wednesday, February 16, \currentYear}
\experiment{Lab Time}
% Place your initial below
NN\\
During the lab time, I continued working on developing a model for the steering system using the Neural Network Time Series toolbox on MATLAB. 

HG\\
During lab I continued analysis of transfer functions to find the one that best represented the brake system. I tried plotting transfer functions with 28 poles and 11 zeros along with transfer functions with 28 poles and 12 zeros to see if the extra zeros would help lower the error to an acceptable range. This only worked for the data set where I used log 2 to build the system and log 3 to validate it. I ended up finding that the transfer function with 28 poles and 10 zeros gave the best fit across data sets, and so I started exporting this transfer function with the goal of re-creating it in Simulink. I chose the transfer function from the data set where log 2 was used to build and log 3 was used to validate because it gave an acceptable best fit and was within the error range. In Simulink I ended up with four transfer function blocks total, since I had four inputs to the system. I'm not sure how to combine all of the transfer functions, so I will need to look into that. 

\labday{Monday, February 21, \currentYear}
\experiment{Lab Time}
% Place your initial below
NN\\
During the lab time, I researched how to develop a neural network using the Neural Network Time Series toolbox on MATLAB and some of the theory behind how neural networks are designed to get a better understanding of what kind of model is being trained and how many delay units and hidden neurons I would need. With this information, I got a working model of the steering subsystem within the desired accuracy requirements. I also met with Dr. Miah to discuss if the same work should be done with the acceleration subsystem since there were some issues with creating a Simulink diagram for the transfer function model. Dr. Miah suggested creating a neural network for the acceleration subsystem.

HG\\
During the lab time I realized that the sampling time for the brake system needs to be set at 33ms, and that I had it set at 34ms. Most of the lab time was then spent re-doing the transfer function analysis for the different data sets. I also researched how to implement inputting a signal to the transfer functions in the Simulink model, along with looking at documentation about the fuzzy logic toolbox. 

Near the start of today's meeting Dr. Miah came in to answer some questions we had. He mentioned that we should make everything in Simulink as automated as possible. We talked about the upcoming progress report, and he mentioned that I should develop a deep learning model for the brake system.  


\experiment{Meeting Minutes}
%Place initial below
HG\\
Since it was President's day, we did not meet with Erik. Instead Dr. Miah, Nick and I met to discuss the project. Dr. Miah encouraged us to look at the report template and to go through the IEEE explorer and look at how people write reports in order to get an idea of how the format should look. We then decided that we would meet with Erik and Joe on Tuesday in order to make up for Monday's meeting. 

\labday{Tuesday, February 22, \currentYear}
\experiment{Meeting Minutes}
% Place your initial below
HG\\
Dr. Miah, Erik, Joe and I met to discuss the progress we had made so far in the senior project. Nick was unable to attend due to an internet outage. We talked about the neural network steering model that Nick had developed, and how we wanted to test it using the HIL bench. There are some concerns that the neural network model might cause trouble with AutonomouStuff's model since it requires the deep learning toolbox, which is something that we need to look into. If needed, the model can be generated into code that can be used by AutonomouStuff. Dr. Miah mentioned that we may need to visit AutonomouStuff's offices a few times, and asked when we could come an test the models. Erik mentioned that he was busy this week, but that the HIL bench is basically ready. Therefore we can just send him the models we have and then set up a time to visit from there. Joe mentioned that they have customer deadlines coming up that may make scheduling the HIL bench more cumbersome going forward. We talked about getting a skeleton diagram, and Erik mentioned that it will look a lot like the other diagrams unless we want the details. I talked a little about what was going on the brake system. For sending the models, AutonomouStuff will be given two sets of models, one using the system identification toolbox and one using the deep learning toolbox. Near the end of the meeting we talked about the progress presentation. In terms of scheduling visits, we will see how the first round of testing goes and then plan from there. 

\labday{Wednesday, February 23, \currentYear}
\experiment{Lab Time}
% Place your initial below
NN\\
During the lab time, I worked on developing a model for the acceleration system using the Neural Network Time Series toolbox on MATLAB. Once I developed a model and exported it to Simulink, I edited the steering and acceleration system neural network models so the sampling time was not set as 1 second, but inherited instead. Also, I had to recreate blocks within the model so that the model can be used without the Deep Learning toolbox. I completed both steering and acceleration system models and they are ready to be sent to Erik to prepare for testing.

HG\\
Once I had re-done all of the transfer functions, I ended up going with the transfer function that had 28 poles and 12 zeros that was generated using the data set where log 2 was used to build and log 3 was used to validate. I ended up choosing this model since it gives an acceptable best fit percentage and because it gives an error within the acceptable error margins. I then exported this transfer function to the MATLAB workspace and set up the transfer function blocks in Simulink. I then started working on generating a neural network model of the brake system. Dr. Miah stopped by at the end of the lab time in order to talk with us about the models. Nick explained that since the neural network model is just a bunch of simple blocks, Erik should be able to run it without having to make use of the deep learning toolbox. We also discussed the need to find a way to get the MISO transfer function system working. We resolved to complete all subsystems using the neural network and then send them to Erik once we know that they work in Simulink without using the deep learning toolbox. Our target to visit AutonomouStuff is the first week of March, and we are also setting a goal to have all of the models done by this Friday. Our progress presentation will be this next Tuesday, and we can use the presentation template we already have. Once we have the three subsystems done, then we will start working on the report. 


\labday{Monday, February 28, \currentYear}
\experiment{Lab Time}
% Place your initial below
NN\\

HG\\
During the lab time I kept working on the brake model I generated using the Neural Network Time Series app. I added a mux block for the voltage inputs, along with sum and scope blocks in order to check the error between the desired and actual pedal positions. I then tested the model on my own computer in order to verify that it could work without the use of the deep learning toolbox. Once I had verified that I could use it on my computer without any problems, I then tested it using constant blocks. This was done to see if the system would generate the desired output. The system was able to generate the desired output and the error was within an acceptable range. Next I tested it using from workspace blocks and the data from the first brake log. Once I did this I found that the error plot had some spikes over 10\%, so I had to retrain the model. It did not take long to retrain the model and find one that gave satisfactory results. I then exported this model to Simulink and changed it to make it able to run without using the deep learning toolbox. I then tested it again using the constant and from workspace blocks, and this time the model passed both tests. 

\labday{Tuesday, March 1, \currentYear}
\experiment{Progress Presentation}
% Place your initial below
HG\\
Nick and I presented the progress we had made so far on the senior project to Dr. Miah, Erik, and Joe. Dr. Miah explained to Nick and I that the transfer function model works well for linear systems, but if you have a highly non-linear system, then the transfer function model will not work. He also said that neural network modeling works best for highly non-linear systems. Once we had given our presentation, we then talked about working to get a date set for us to visit in order to test our models and collect more data. Dr. Miah mentioned that we have eight weeks exactly left to complete the project. We want to go through the first iteration of testing as soon as possible. Joe said that he would look into moving things around in order to make this possible. In the meantime we will work on the report. Erik said that he has started working on getting the HIL harness ready. He may be getting help with this, but it depends on who's available. Near the end of the meeting Joe mentioned that we should get this completed within the next two weeks. 

\labday{Wednesday, March 2, \currentYear}
\experiment{Lab Time}
% Place your initial below
NN\\


HG\\
During the lab time I first worked on exporting the neural network brake model to MATLAB 2019b. This was done to ensure that it could be used by AutonomouStuff. I then worked on the report, mainly focusing on the abstract and the explanation of why we are using neural network to create models. 


\labday{Monday, March 7, \currentYear}
\experiment{Lab Time}
% Place your initial below
HG\\
During the lab time I worked on the report. I added and formatted text about why we are using the Neural Network in the main project work section. I then reproduced the error and actual pedal position plots from the brake subsystem using the brake Neural Network model. I did this in order to create plots that could be used in the report. I then generated the same plots but I used constant blocks as the model inputs. I worked on the brake system section of the report under validation and testing, and added some plots. After the meeting Dr. Miah asked me to let him know when I pushed all of the changes to the repository so he could start giving us feedback on the report. 

\experiment{Meeting Minutes}
% Place your initial below
HG\\
Today Dr. Miah, Erik, and I met near the end of the lab time. Our main topic of discussion was setting up a date for testing. Erik mentioned that he had received our updated models and passed them onto the HIL engineer. He is going to check with him to make sure that they are working. Erik also mentioned that this Wednesday they are having a meeting to discuss the project and see if anyone else from the company would be able to help us. He will let us know what happens after the meeting. 


\labday{Wednesday, March 9, \currentYear}
\experiment{Lab Time}
% Place your initial below
NN\\



HG\\
During the lab time I worked on adding more text and figures to the brake system subsection in the validation and testing section of the report. I also tried to change the text above the arrows in the acceleration block diagram to fit, but couldn't do this without making the arrows longer. After searching the internet, I still wasn't able to solve the problem, and Nick mentioned that he had run into the same thing. I then emailed Dr. Miah about this for advice. I added some keywords to the beginning of the report, along with re-doing the from workspace error plot for the brake system in order to ensure it was the same size as the other plots. Near the end of lab I started working on the modeling vehicle subsystem section, but haven't added the text to the report yet. 

\labday{Monday, March 21, \currentYear}
\experiment{Lab Time}
% Place your initial below
NN\\
During the lab, I worked on gathering information on neural network theory to add details to the report. Hannah and I talked with Dr. Miah to get a clearer understanding of the structure of the report. Dr. Miah helped us reorganize the layout of the report so it makes more sense.

HG\\
Today during the lab I worked on reading through what we had so far in the report and made some minor edits. I also started looking at the final presentation and made some changes to that. Dr. Miah stopped by before the weekly meeting with AutonomouStuff to answer some questions we had about the report and the presentation. For the presentation, he explained that it will essentially be the same as the final report. We will need to include a deep learning section, and we decided that instead of mentioning deep learning we will call it Neural Network modeling in our report, and then work in deep learning. 

\experiment{Meeting Minutes}
HG\\ Near the end of the lab time Nick, Dr. Miah, and I met via Zoom with Sai Chittemsetty and Matt Goben from AutonomouStuff. Dr. Miah let them know what we had worked on for this project and where we were currently at. Since we are not going to have the HIL bench, we needed to figure out another way to test the models. Sai mentioned that the company is at a busy time right now, and that he looked at the MATLAB files we sent. After looking at those files he determined that it wouldn't take much time to change the software and be able to do some basic testing. It will be more like a demonstration, with an open loop instead of HIL. This will allow Nick and I to see the real life system and the output on the scope, even though we won't have feedback to see the system performance. Dr. Miah mentioned that we want to finish all testing by April 15th, since that is when we will start working on the final report and presentation. After discussing a few different dates to meet we eventually settled on having a virtual meeting on March 30th at 2pm, and two in-person meetings on April 1st and 4th at 2pm. These shorter sessions will allow Nick and I time to better understand how the whole process works. In the meantime we will update the report. We will have to coordinate with Sai to send him the rest of the MATLAB files as necessary. 
\linebreak
\linebreak
Afterwards Nick and I met with Dr. Miah to discuss some questions about the report. Dr. Miah let us know how to order the authors and about the project work,  section. He reformatted some sections of the paper, and we will work on that this coming Wednesday. 

\labday{Wednesday, March 23, \currentYear}
\experiment{Lab Time}
% Place your initial below
NN\\
I worked on adding information to the neural network modeling, modeling vehicle subsystems, and validation & testings sections of the report.

HG\\
I worked on adding information to the Neural Network modeling, modeling vehicle subsystems, and validation and testing sections of the report. 

%--------------------------------------------------------------------------------------------

\end{addmargin}

%----------------------------------------------------------------------------------------
%    BIBLIOGRAPHY
%----------------------------------------------------------------------------------------


\bibliographystyle{plain}
\bibliography{bib/references.bib}


% \begin{thebibliography}{9}



% \end{thebibliography}

%----------------------------------------------------------------------------------------

\end{document}


%%% Local Variables:
%%% mode: latex
%%% TeX-master: t
%%% End:
